\uuid{md6a}
\titre{Primitives usuelles}
\auteur{hanani ; mzari}
\organisation{exo7}
\licence{CC BY-NC-SA 4.0}
\datecreate{2020-01-20}
\niveau{L1}
\module{Analyse}
\chapitre{Calcul d'intégrales}
\sousChapitre{Primitives diverses}
\difficulte{2}
\theme{Primitives, primitives usuelles}
\exo7qcmid{lille-2.primitives.011}

\contenu{
\begin{question}
\qclassification{127.11}
\qtags{motcle=primitives usuelles}
Parmi les affirmations suivantes, cocher celles qui sont vraies :
\begin{answers}
\good{Sur $]-1,+\infty[$, on a : $\displaystyle \int \frac{x^2\mathrm{d}x}{1+x^3}=\frac{1}{3}\ln (1+x^3)+k$, où $k\in \Rr$.}
\bad{Sur $]-1,+\infty[$, on a : $\displaystyle \int \frac{x^2\mathrm{d}x}{1+x^3}=\frac{x^3}{3(1+x^3)}+k$, où $k\in \Rr$.}
\bad{Sur $]1,+\infty[$, on a : $\displaystyle \int \sqrt{x-1}\mathrm{d}x=\frac{1}{\sqrt{x-1}}+k$, où $k\in \Rr$.}
\good{Sur $]1,+\infty[$, on a : $\displaystyle \int \sqrt{x-1}\mathrm{d}x=\frac{2}{3}(x-1)^{3/2}+k$, où $k\in \Rr$.}
\end{answers}
\vskip3mm
\begin{explanations}
Avec $u=1+x^3$, on a : $\mathrm{d}u=3x^2\mathrm{d}x$ et
$$\int \frac{x^2\mathrm{d}x}{1+x^3}=\frac{1}{3}\int \frac{\mathrm{d}u}{u}=\frac{1}{3}\ln |u|+k=\frac{1}{3}\ln (1+x^3)+k,\; k\in \Rr.$$
De m\^eme, avec $u=x-1$, on a $\mathrm{d}u=\mathrm{d}x$ et
$$\int \sqrt{x-1}\mathrm{d}x=\int \sqrt{u}\mathrm{d}u=\frac{2}{3}u^{3/2}+k=\frac{2}{3}(x-1)^{3/2}+k,\; k\in \Rr.$$
\end{explanations}
\end{question}
}
